\epigraph{herramientas y tecnologías}{

    La implementación de la solución propuesta requiere la selección cuidadosa de tecnologías, 
    herramientas y \textit{frameworks} que permitan materializar de manera eficiente los modelos 
    curriculares, soportar reglas de variabilidad y garantizar una experiencia de desarrollo fluida. 
    En este epígrafe se describen los componentes tecnológicos que conformarán la plataforma, 
    justificando su elección en función de Solo FastAPIs como escalabilidad, mantenibilidad, 
    interoperabilidad y robustez. Asimismo, se analizan los entornos de desarrollo, 
    lenguajes de programación, sistemas de persistencia y bibliotecas especializadas que 
    facilitan la representación formal de restricciones curriculares y la validación 
    automática de configuraciones. Esta caracterización permite comprender la base 
    tecnológica que sustenta la solución y cómo cada herramienta contribuye al logro de 
    sus objetivos funcionales y no funcionales.

    \subepigraph{Lenguaje de Programación}{

        Python será el lenguaje principal utilizado para el desarrollo del \textit{backend} de la 
        solución propuesta. Su elección se fundamenta en su madurez tecnológica, su amplio 
        ecosistema de librerías y la claridad sintáctica que facilita la implementación de 
        modelos complejos manteniendo niveles altos de legibilidad y mantenibilidad del código 
        \citep{van1995python, lutz2013learning}. Python es ampliamente utilizado en el ámbito 
        académico y profesional, lo que garantiza disponibilidad de documentación, ejemplos de 
        buenas prácticas y una sólida comunidad de soporte \citep{ray2014large}.

        Una de las principales fortalezas de Python radica en su capacidad de integración 
        con múltiples paradigmas de desarrollo, permitiendo tanto programación estructurada 
        como orientada a objetos y funcional \citep{lutz2013learning}. Esto resulta particularmente 
        útil para el diseño modular del sistema, la definición formal de \ac{FM} y la implementación 
        de validaciones de reglas en motores lógicos.

        La representación estructurada de modelos de características \ac{FM} requiere un lenguaje 
        capaz de expresar jerarquías, relaciones lógicas, restricciones transversales y 
        ardinalidades con claridad y sin ambigüedad. Python ofrece capacidades idóneas para 
        este propósito debido a su flexibilidad en estructuras de datos y su riqueza semántica 
        para definir invariantes y reglas de negocio.

        En primer lugar, Python permite modelar FM de forma natural mediante diccionarios, 
        listas y clases, manteniendo la semántica jerárquica del dominio sin necesidad de 
        sintaxis complejas ni compilación intermedia. Esto facilita representar nodos, ramas, 
        grupos OR/XOR, atributos y relaciones cross-tree como estructuras anidadas. Gracias a 
        esta expresividad, resulta posible describir distintos niveles del modelo —desde 
        características obligatorias y opcionales hasta restricciones de cardinalidad— de manera 
        directa y legible, lo que favorece tanto su mantenimiento como su inspección.

        Además, Python ofrece un ecosistema amplio de herramientas para la validación 
        programática de configuraciones. El backend puede implementar algoritmos de selección y 
        verificación que comprueben automáticamente consistencia lógica, satisfacibilidad de 
        prerrequisitos, exclusiones y cumplimiento de constraints. La sintaxis de Python resulta 
        especialmente adecuada para traducir expresiones lógicas como requires, excludes o 
        IMPLIES, ya sea mediante evaluaciones directas o a través de módulos que permiten 
        construir árboles lógicos, satisfacibilidad booleana o incluso resolución por backtracking.

        Otro aspecto clave es que Python permite separar con claridad el modelo 
        (estructura estática) de las configuraciones (instancias dinámicas). A través de 
        funciones o clases especializadas, el sistema puede generar configuraciones válidas a 
        partir de elecciones de usuario, aplicar reglas de minimización/maximización de 
        cardinalidades, determinar incompatibilidades, detectar violaciones y ofrecer 
        soluciones alternativas, todo ello sin necesidad de traducir el modelo a lenguajes 
        externos. Adicionalmente, la flexibilidad del lenguaje, junto con el respaldo de su 
        amplio ecosistema de librerías, facilita la construcción de algoritmos que no solo 
        derivan configuraciones de forma eficiente, sino que también las verifican mediante 
        evaluaciones lógicas, análisis de dependencias y resolución de conflictos. Este mismo 
        soporte permite implementar mecanismos de versionado y snapshots del modelo, conservando 
        estados históricos y garantizando que la evolución del FM no invalide configuraciones 
        previas, preservando así la consistencia y trazabilidad del sistema en el tiempo.

        Asimismo, Python ofrece un ecosistema robusto para el desarrollo de aplicaciones web 
        a través de frameworks como FastAPI, que permiten construir servicios RESTful eficientes 
        y escalables, lo cual facilita exponer un motor de validación y generación de configuraciones 
        como servicio \citep{ramirez2019fastapi}. A esto se suma la compatibilidad con herramientas 
        modernas de persistencia de datos como SQLAlchemy y pydantic, lo que facilita la representación 
        estructurada de información curricular y su validación automática. Esto habilita una arquitectura 
        limpia en la que las representaciones de \ac{FM}, los algoritmos de derivación y las reglas 
        del dominio educativo quedan encapsulados y accesibles desde otros componentes de la 
        plataforma (cliente, bases de datos, módulos administrativos, etc.).

        En conjunto, estas características hacen de Python un lenguaje altamente apropiado 
        tanto para la representación formal de modelos de características como para la generación 
        y validación automatizada de configuraciones dentro del sistema propuesto \citep{ray2014large}.
        Por su equilibrio entre simplicidad expresiva y potencia funcional, Python constituye 
        una base tecnológica idónea para el \textit{backend} del sistema, permitiendo acelerar el 
        desarrollo sin sacrificar calidad, extensibilidad ni alineación con los requisitos 
        funcionales de la solución.

    }

    \subepigraph{Framework de Desarrollo}{
        
        FastAPI será el framework \textit{backend} adoptado para la implementación de la plataforma, 
        debido a su equilibrio entre simplicidad de desarrollo, alto rendimiento y robustez 
        arquitectónica \citep{ramirez2019fastapi, eugene2023fastapi}. Su diseño se basa en los 
        principios de asincronía nativa y validación automática de datos, permitiendo construir 
        APIs eficientes, escalables y seguras, características fundamentales para un sistema que 
        debe procesar configuraciones de modelos curriculares, validar restricciones y ofrecer 
        respuestas consistentes en tiempo real.

        Un factor determinante en la selección de FastAPI es su rendimiento comparable a 
        implementaciones escritas en Go y Node.js, gracias a su construcción sobre Starlette 
        (para la capa asíncrona) y Pydantic (para validación y serialización) \citep{pesce2021benchmark}. 
        Esta combinación permite manejar múltiples peticiones concurrentes, realizar verificaciones 
        de reglas e interacciones dinámicas con el motor de configuraciones sin degradación perceptible 
        del servicio, incluso bajo cargas elevadas.

        La validación de datos es un aspecto crítico del sistema propuesto, ya que el backend 
        debe recibir elecciones de características, evaluarlas frente a cardinalidades y 
        restricciones lógicas, y retornar configuraciones válidas o alternativas corregidas. 
        En este sentido, FastAPI destaca por integrar modelos de datos declarativos mediante 
        Pydantic, lo que garantiza que cada solicitud sea verificada estructural y semánticamente 
        antes de su procesamiento \citep{ramirez2019fastapi}. Esto no solo previene errores de 
        ejecución, sino que contribuye a la robustez del motor de configuraciones, evitando 
        inconsistencias desde la capa de entrada.

        Un argumento especialmente relevante para seleccionar FastAPI es la naturaleza 
        intensiva en validaciones y operaciones I/O que caracteriza al sistema propuesto. 
        La plataforma no solo debe recibir configuraciones y verificar reglas, sino también 
        consultar persistencia, aplicar evaluaciones lógicas, generar respuestas alternativas y 
        registrar versiones de modelo, todo ello de manera concurrente y potencialmente bajo 
        múltiples usuarios. Debido a su arquitectura asíncrona nativa basada en async/await, 
        FastAPI permite manejar este patrón de carga sin bloqueo, maximizando la eficiencia en 
        tareas I/O y reduciendo latencias \citep{eugene2023fastapi, pesce2021benchmark}. Esto se 
        traduce en una capacidad real de escalar, sin incurrir en cuellos de botella característicos 
        de frameworks síncronos, lo que resulta esencial en un sistema donde las consultas y 
        validaciones son frecuentes, complejas y potencialmente concurrentes.

        Adicionalmente, FastAPI presenta una curva de aprendizaje notablemente baja, 
        permitiendo extremo enfoque en la lógica de negocio sin sacrificar mantenibilidad. 
        Su documentación generada automáticamente (OpenAPI/Swagger) mejora la trazabilidad 
        del sistema y facilita la colaboración entre desarrolladores y expertos curriculares, 
        quienes pueden visualizar endpoints, contratos de datos y respuestas sin requerir 
        interacción directa con el código \citep{ramirez2019fastapi}.

        Desde el punto de vista arquitectónico, FastAPI favorece un desarrollo altamente 
        modular, encajando de manera natural con principios de separación de responsabilidades y 
        diseño por componentes \citep{richards2015software}. Esta modularidad resulta especialmente 
        útil para organizar servicios de modelado de FM, motor de restricciones, evaluadores de 
        configuraciones, repositorios, controladores y almacenes de versión.

        Finalmente, la compatibilidad de FastAPI con controladores de bases de datos SQL, 
        junto con herramientas de orquestación (Docker) y despliegue CI/CD, permite 
        que la solución crezca de forma controlada desde prototipos académicos hasta implementaciones 
        institucionales. En conjunto, su velocidad, asincronía, simplicidad expresiva, capacidad de
        validación y flexibilidad lo convierten en la opción más adecuada para el sistema curricular 
        propuesto \citep{eugene2023fastapi}.
        
    }
    
    \subepigraph{Base de Datos y Gestor de Base de Datos}{
        
        La persistencia de los modelos curriculares, configuraciones generadas, 
        versiones históricas y metadatos asociados requiere de un sistema robusto, 
        seguro y con garantías de consistencia transaccional. En este proyecto se 
        selecciona PostgreSQL como gestor de base de datos relacional, debido a su madurez, 
        fiabilidad y su capacidad de manejar estructuras complejas con fuertes requisitos 
        de integridad referencial. PostgreSQL destaca por su cumplimiento estricto del 
        estándar SQL, su modelo ACID, y su eficiente gestión de transacciones, elementos 
        imprescindibles para representar relaciones jerárquicas entre características, 
        restricciones cross-tree, y snapshots versionados del modelo sin comprometer la 
        coherencia global del sistema.

        Además, PostgreSQL ofrece capacidades avanzadas como índices compuestos, consultas 
        optimizadas, vistas materializadas, triggers y procedimientos almacenados, que 
        resultan útiles para procesar datos curriculares con alto grado de estructuración. 
        El motor también soporta tipos de datos JSON y JSONB, permitiendo almacenar 
        representaciones semiestructuradas de configuraciones y estados transitorios, por 
        lo que proporciona flexibilidad para almacenar tanto esquemas estáticos como 
        snapshots dinámicos relacionados con la evolución del Modelo de Características.

        Adicionalmente, PostgreSQL se integra de forma natural con SQLAlchemy como ORM 
        dentro del backend desarrollado en FastAPI, facilitando el diseño declarativo de 
        entidades, la gestión del ciclo de persistencia, la validación de relaciones y la 
        migración segura del esquema. Esta integración favorece una arquitectura limpia y 
        escalable, desacoplando las capas de datos, servicios y lógica de negocio.

        Para la administración visual del motor se empleará pgAdmin, interfaz ampliamente 
        utilizada para la gestión de bases de datos PostgreSQL. Mediante esta herramienta 
        es posible inspeccionar tablas, ejecutar consultas SQL, visualizar índices y 
        constraints, supervisar cargas de trabajo y realizar respaldos de forma controlada. 
        En el contexto de la solución propuesta, pgAdmin permite auditar el estado de los 
        modelos almacenados, verificar reglas impuestas, revisar configuraciones generadas 
        y garantizar la trazabilidad del sistema.

        En conjunto, la elección de PostgreSQL como sistema de persistencia y pgAdmin 
        como herramienta administrativa proporciona una base confiable y alineada 
        con las necesidades de integridad, consistencia y trazabilidad del sistema 
        curricular modelado.
            
    }

    \subepigraph{Contenerización y Orquestación}{

        Para garantizar la portabilidad, reproducibilidad y facilidad de despliegue del 
        sistema propuesto, se adoptará Docker como tecnología base de contenedorización. 
        Docker permite encapsular los servicios de la aplicación —incluyendo el backend 
        FastAPI, la base de datos PostgreSQL, la herramienta administrativa pgAdmin y
        otros servicios de utilidad o futuros módulos asociados— dentro de entornos aislados 
        y autocontenidos. Esto elimina problemas de dependencias, divergencias de entorno
        y configuraciones inconsistentes entre despliegues locales y productivos.
        La capacidad de reproducir una imagen de forma determinística contribuye directamente 
        a la robustez y mantenibilidad de la solución.

        La contenedorización resulta especialmente útil para aplicaciones multicapa como 
        la propuesta, donde distintos servicios deben colaborar manteniendo canales de 
        comunicación estables. Docker permite garantizar uniformidad tanto en las 
        dependencias del backend, como versiones de librerías, configuración del servidor 
        de aplicación y parámetros del engine de base de datos, disminuyendo errores 
        derivados de diferencias de entorno y facilitando diagnósticos.

        Sobre Docker se construye Docker Compose, herramienta de orquestación ligera que 
        permite describir, mediante un único archivo declarativo, toda la topología de 
        servicios del sistema. Esto incluye la configuración de redes internas, gestión 
        de persistencia en volúmenes, variables de entorno, reglas de dependencias entre 
        contenedores, puntos de exposición de puertos, así como políticas de reinicio. 
        En el contexto del sistema de modelado curricular, Docker Compose facilita levantar 
        el entorno completo de desarrollo o producción con un solo comando, garantizando 
        escalabilidad y reduciendo tiempos de despliegue o recuperación ante fallos.

        Un beneficio adicional de esta estrategia es su contribución a las prácticas de 
        Integración Continua y Despliegue Continuo (CI/CD). Gracias a la naturaleza 
        declarativa y reproducible de los contenedores, el sistema puede desplegarse 
        automáticamente mediante pipelines controlados sin intervención manual, asegurando 
        que las versiones del modelo curricular, las configuraciones generadas y los 
        snapshots almacenados sean accesibles en entornos iguales o equivalentes entre sí.

        De esta forma, el uso conjunto de Docker y Docker Compose aporta portabilidad,
        eficiencia en despliegue, coherencia entre entornos y una base técnica sólida para 
        escalar la solución futura en producción.
            
    }
    
    \subepigraph{Almacenamiento en Memoria, Caching y Mensajería}{

        Redis será incorporado como componente clave para el almacenamiento temporal 
        y procesamiento eficiente de información crítica durante las fases de 
        generación y validación de configuraciones curriculares. Su naturaleza de 
        base de datos en memoria permite alcanzar tiempos de respuesta extremadamente 
        bajos, lo cual resulta particularmente valioso en un sistema donde las 
        operaciones de verificación de restricciones, derivación de configuraciones 
        válidas y resolución de incompatibilidades pueden ejecutarse de forma intensiva 
        y repetitiva.

        A diferencia de un motor relacional tradicional, Redis está optimizado 
        para operaciones clave–valor y estructuras avanzadas como listas, conjuntos, 
        mapas hash y sorted sets. Estas estructuras permiten modelar transitoriamente 
        elementos como:
        \begin{itemize}
            \item configuraciones candidatas generadas por el usuario,
            
            \item listas de incompatibilidades detectadas,
            
            \item descriptores internos del Feature Model,
            
            \item estados de sesión durante la exploración de alternativas,
            
            \item cachés de resultados de validaciones repetitivas.
        \end{itemize}

        Este almacenamiento en memoria reduce significativamente la latencia, lo que 
        contribuye a experiencias interactivas y a la capacidad del sistema de 
        responder en tiempo real a decisiones del usuario (por ejemplo, activar 
        o desactivar características, consultar prerequisitos, etc.).

        Otro aspecto clave es que Redis permite persistencia opcional mediante 
        snapshots (RDB), lo que lo convierte en un aliado directo del soporte de 
        versionado de modelos y configuraciones. Es decir, puede almacenar estados 
        intermedios del proceso de modelado curricular sin interferir con la base de 
        datos principal, permitiendo restaurar fácilmente versiones anteriores o 
        explorar ramas alternativas del mismo modelo.

        Además, Redis implementa un sistema de Time-To-Live (TTL), lo que permite 
        gestionar la caducidad automática de configuraciones temporales, sesiones y 
        cálculos ya procesados. Esta funcionalidad contribuye a la integridad global 
        del sistema y evita acumulación de datos residual.

        Finalmente, una característica determinante para su elección es su excelente 
        integración con FastAPI, permitiendo aprovechar plenamente la naturaleza 
        asíncrona del backend. Gracias a librerías como aioredis, las operaciones de 
        lectura/escritura se realizan sin bloquear el event loop, lo que maximiza la 
        concurrencia cuando múltiples usuarios exploran configuraciones simultáneamente.

        En síntesis, Redis aporta velocidad, soporte a cálculos temporales complejos, 
        persistencia selectiva, alta concurrencia y facilidad de integración con el 
        framework propuesto, consolidándose como la pieza ideal para manejar la lógica 
        dinámica del sistema.
            
    }
    
    \subepigraph{Procesamiento Asíncrono}{

        En la arquitectura propuesta, Celery desempeñará un rol esencial como 
        orquestador de tareas asíncronas, permitiendo que operaciones complejas, 
        lentas o intensivas en cómputo se ejecuten fuera del ciclo de solicitud-respuesta 
        del API, mejorando sustancialmente la experiencia del usuario y la disponibilidad 
        del sistema.

        Dado que el proceso de generación y validación de configuraciones curriculares 
        incluye fases donde se deben:
        \begin{itemize}
            \item Recorrer árboles completos del Feature Model,
            
            \item Validar múltiples restricciones cruzadas,
            
            \item Propagar efectos lógicos por dependencias,
            
            \item Explorar configuraciones alternativas válidas,
            
            \item Verificar versiones y retrocompatibilidad,
        \end{itemize}

        Resulta natural que estas tareas no se ejecuten directamente en una petición HTTP, 
        ya que podrían bloquear el servidor o provocar tiempos de respuesta superiores 
        a lo aceptable. Celery permite delegar estos procesos a “workers” paralelos 
        que pueden ejecutarse en segundo plano, liberando a FastAPI para seguir 
        atendiendo solicitudes sin congestión.

        La integración de Celery con Redis (como message broker) permite que las tareas 
        se publiquen, distribuyan, encolen, ejecuten y documenten de forma confiable. 
        Asimismo, Redis funge como backend de resultados, permitiendo al sistema:
        \begin{itemize}
            \item Consultar el estado de una ejecución,
            
            \item Recuperar resultados parciales o finales,
            
            \item Cancelar o reprogramar tareas,
            
            \item Mostrar progresos al usuario.
        \end{itemize}

        Otra ventaja crítica es que Celery soporta escalamiento horizontal, lo cual es 
        relevante para tu sistema, en el cual múltiples usuarios pueden realizar 
        evaluaciones simultáneamente sobre modelos distintos o versiones diferentes. 
        Con más carga, simplemente se añaden más workers, sin necesidad de cambios 
        estructurales.

        Desde el punto de vista de robustez, Celery ofrece herramientas para:
        \begin{itemize}
            \item Reintentos automáticos ante fallos,
            
            \item Límite de concurrencia,
            
            \item Priorización de colas,
            
            \item Periodicidad (con Celery Beat) para ejecutar tareas programadas.
        \end{itemize}

        Esto permite automatizar procesos como:
        \begin{itemize}
            \item Generación periódica de snapshots de Feature Models,
            
            \item Validaciones nocturnas de consistencia,
            
            \item Limpieza de configuraciones expiradas,
            
            \item Backups de estados intermedios.
        \end{itemize}

        Finalmente, Celery es altamente compatible con FastAPI y Python, manteniendo 
        coherencia tecnológica con la decisión de usar un stack asíncrono y un modelo 
        de validación compleja. Su adopción garantiza que la capa lógica pesada no 
        degrade la utilidad interactiva del sistema, sino que se ejecute de forma 
        paralela, programada y resistente. La siguiente tabla \ref{tab:fastAPIvsFastAPIconCelery} 
        respalda la justificación de por qué adicionar celery al sistema.

        \begin{longtable}{|p{5cm}|p{5cm}|p{5cm}|}
            \caption{Tabla Comparativa: FastAPI vs FastAPI + Celery \label{tab:fastAPIvsFastAPIconCelery}}\\
            \hline
            \rowcolor[HTML]{CBCEFB} 
            \textbf{Criterio} & \textbf{Solo FastAPI} & \textbf{FastAPI + Celery} \\
            \hline
            \endfirsthead
            
            \multicolumn{3}{c}
	        {{\bfseries \tablename\ \thetable{} -- Continuación de la página anterior}} \\
            
            \hline
            \rowcolor[HTML]{CBCEFB} 
            \textbf{Criterio} & \textbf{Solo FastAPI} & \textbf{FastAPI + Celery} \\
            \hline
            \endhead
            
            \hline
            \multicolumn{3}{|r|}{{Continúa en la siguiente página}} \\
            \endfoot
            
            \hline
            \endlastfoot
            
            % Cuerpo de la tabla
            Manejo de operaciones pesadas & Bloquea el servidor si la tarea tarda mucho & Delegación a workers externos sin bloqueo \\ 
            \hline
            \rowcolor[HTML]{EDEBF1}
            Experiencia del usuario & Tiempos de respuesta largos para validaciones & Respuestas inmediatas + seguimiento del progreso \\ 
            \hline
            Escalabilidad & Limitada al número de workers del servidor & Escalamiento horizontal agregando workers Celery \\ 
            \hline
            \rowcolor[HTML]{EDEBF1}
            Naturaleza de carga & Orientada a I/O y sincronización & Orientada a cómputo intensivo y tareas diferidas \\ 
            \hline
            Validación de modelos complejos & Puede agotar recursos del servidor & Optimización mediante workers paralelos \\ 
            \hline
            \rowcolor[HTML]{EDEBF1}
            Generación de snapshots o versiones & Bajo rendimiento si se ejecuta en el request & Tareas programadas y automáticas (Celery Beat) \\ 
            \hline
            Robustez ante fallos & Fallos interrumpen solicitudes activas & Reintentos automáticos e independientes \\ 
            \hline
            \rowcolor[HTML]{EDEBF1}
            Concurrencia & Limitada & Altamente concurrente \\ 
            \hline
            Distribución de carga & Manual, difícil de gestionar & Balance natural mediante colas y brokers \\ 
            \hline
            \rowcolor[HTML]{EDEBF1}
            Estado de procesos largos & No persistente, no rastreable & Persistente, rastreable, recuperable \\ 
            \hline
            Integración con Redis & Opcional & Ideal, Redis como broker y backend \\ 
            \hline
            \rowcolor[HTML]{EDEBF1}
            Mantenimiento & Monolítico & Arquitectura desacoplada y modular \\ 
            \hline

        \end{longtable}

        Esta comparación evidencia que mientras FastAPI permite construir servicios web 
        altamente eficientes para operaciones sin bloqueo y validaciones ligeras, la 
        incorporación de Celery habilita un procesamiento robusto y distribuido ideal 
        para la naturaleza intensiva de la manipulación y validación de Feature Models, 
        generando configuraciones, snapshots y cálculos alternativos sin afectar la 
        disponibilidad del sistema.

            
    }
    
    \subepigraph{Algoritmos, Motor de Reglas y Validación}{
        algoritmos geneticos, IA, SymPy y NetworkX
    }
    
    \subepigraph{Gestor de dependencias}{
        poetry y uv
    }
    
    \subepigraph{Herramienta para el control de versiones}{
        GIT y GITHUB
    
    }
    
    \subepigraph{Herramienta y lenguaje de modelado}{
        uml y visual paradigm
    }
    
    \subepigraph{Herramientas para diseñar y ejecutar pruebas}{
        postman y pruebas unitarias
    }
    
    \subepigraph{Herramientas para diseñar y ejecutar pruebas}{
        Structlog Sentry
    }
    
    \subepigraph{Entorno Integrado de Desarrollo}{

        Los Entornos Integrados de Desarrollo (IDE) son las herramientas que 
        todo desarrollador debe tener a mano. Permiten editar y gestionar
        código fuente en diversos lenguajes de programación y ofrecen
        múltiples herramientas para facilitar el trabajo y aumentar la productividad, pues
        permiten depurar, resaltar sintaxis, autocompletado, integración con control de versiones
        y la refactorización. Por lo general cada IDE está asociado a un único 
        lenguaje de programación o a un conjunto pequeño de lenguajes, cuyo factor 
        provoca que la gran mayoría de desarrolladores opten por otras alternativas, 
        ya que en proyectos donde el uso de múltiples lenguajes es evidente, 
        puede ser un obstáculo \citep{ibm_IDE}.

        Se escogió como entorno de desarrollo Visual Studio Code \citep{vscode2025}, cuya 
        herramienta es no convencional pero muy usada por muchos programadores, 
        siendo señalado como el editor de código más popular entre proyectos de código abierto 
        (\textit{open source}) \citep{ramel2021VSCodePopularity}.
        El creador del editor de código multiplataforma, VSCode, es Microsoft. Este software es
        gratuito y de código abierto. Incluye soporte para la depuración, control
        integrado de Git, resaltado de sintaxis, finalización inteligente de código,
        fragmentos y refactorización de código. Además, 
        su nivel de personalización permite adaptar el entorno a las preferencias de cada 
        programador mediante temas, atajos y configuraciones avanzadas \cite{vscode_docs2025}.

        VSCode presenta extensiones importantísimas que permite adaptar el
        entorno de desarrollo para una enorme infinidad de tareas, tecnologías y lenguajes. 
        Este conjunto, prácticamente ilimitado, es totalmente accesible
        dentro del propio editor. Tiene soporte integrado para aplicaciones web;
        por lo tanto, se puede llevar a cabo en este entorno de trabajo la
        codificación de la aplicación a desarrollar con las tecnologías Fast API y
        Python. VSCode es considerado por brindar eficiencia y agilidad en la
        programación, pues funciona muy bien en incluso equipos con recursos
        limitados; además, los desarrolladores lo aprecian porque su interfaz de
        usuario es muy intuitiva y permite comenzar, prácticamente, a trabajar sin
        necesidad de alguna experiencia \cite{vscode2025, vscode_docs2025}.

    }
}
